%! Vector-Valued Functions — Unit 4, Lesson 4.9 — Lesson Plan
\documentclass[10pt]{article}
\usepackage{precalculus-article}
\usepackage{precalculus-boxes}
\usepackage{graphicx}
\graphicspath{{images/}}
\newcommand{\TallMath}[1]{$\displaystyle #1\rule[-1.4em]{0pt}{3.2em}$}
\newcommand{\CourseName}{AP Precalculus}
\newcommand{\SchoolYear}{2026--2027}
\newcommand{\MeetingLength}{55 minutes}
\newcommand{\UnitNumberName}{Unit 4: Functions Involving Parameters, Vectors, and Matrices \quad}
\newcommand{\LessonNumberName}{Lesson 4.9: Vector-Valued Functions}

\begin{document}
\begin{center}
    {\Large\bfseries \CourseName: \SchoolYear \\
    {\small\bfseries \UnitNumberName \LessonNumberName}}
\end{center}

\begin{tcolorbox}[colback=sky, colframe=navy, boxrule=0.9pt, arc=2mm,
    left=2mm, right=2mm, top=1.5mm, bottom=1.5mm]
    \textbf{Primary Objective:} Students express parametric planar motion as
    vector-valued position and velocity functions, and interpret the magnitude
    of these vectors in context.
    \hfill \textit{(Big Idea: Functions)}
\end{tcolorbox}

\renewcommand{\arraystretch}{1.6}
\begin{skillbox}[Priority Ideas \& Skills]{goldbox}
\begin{tabularx}{\linewidth}{@{}p{2.6cm}X@{}}
\textbf{AP Skill 3.C} & Support conclusions with logical rationale or appropriate data.\\
\textbf{EK 4.9.A.1}  & A parametric function $f(t)=(x(t),y(t))$ can be expressed as the
  vector-valued function $\vec{p}(t)=\langle x(t),\,y(t)\rangle$. The magnitude
  $\|\vec{p}(t)\|=\sqrt{[x(t)]^2+[y(t)]^2}$ gives the distance of the particle from
  the origin at time $t$.\\
\textbf{EK 4.9.A.2}  & The velocity vector $\vec{v}(t)=\langle x'(t),\,y'(t)\rangle$
  encodes direction: the sign of $x'(t)$ indicates right/left motion; the sign of
  $y'(t)$ indicates up/down motion. The magnitude $\|\vec{v}(t)\|$ is the speed.
\end{tabularx}
\end{skillbox}

\begin{skillbox}[{Vocabulary, Concepts, \& Theorems}]{greenbox}
\begin{itemize}
  \item \textbf{Vector-valued function} --- a function whose output is a vector,
        $\vec{p}(t)=\langle x(t),\,y(t)\rangle=x(t)\,\hat{\imath}+y(t)\,\hat{\jmath}$
  \item \textbf{Position vector} --- $\vec{p}(t)$; its magnitude gives distance from origin
  \item \textbf{Velocity vector} --- $\vec{v}(t)=\langle x'(t),\,y'(t)\rangle$;
        components encode direction of motion
  \item \textbf{Speed} --- $\|\vec{v}(t)\|=\sqrt{[x'(t)]^2+[y'(t)]^2}$, the
        magnitude of the velocity vector
\end{itemize}
\end{skillbox}

\begin{skillbox}[Activate Prior Knowledge \& Spiral Review]{sky}
\begin{tabularx}{\linewidth}{@{}X@{}}
\textbf{Lesson 4.1--4.2 (Parametric Functions \& Planar Motion):}
Evaluating $f(t)=(x(t),y(t))$ at a given $t$; interpreting $x(t)$ and $y(t)$
as coordinates; direction of motion from sign of average rates of change.\\[3pt]
\textbf{Lesson 4.8 (Vectors):}
Component form $\langle a,\,b\rangle$; vector magnitude
$\|\langle a,b\rangle\|=\sqrt{a^2+b^2}$; unit vectors $\hat{\imath}$, $\hat{\jmath}$.
\end{tabularx}
\end{skillbox}

\begin{skillbox}[Hook]{sky}
A drone's position at time $t$ seconds is given by $x(t)=3t$ and $y(t)=t^{2}-4$
(distances in meters, launch at the origin).
\begin{enumerate}
  \item How far from the launch point is the drone at $t=2$?
  \item Is the drone moving up or down at $t=2$?  Left or right?
  \item Can a single mathematical object capture both the position \emph{and} the
        motion at once?  What might it look like?
\end{enumerate}
\textit{Goal: motivate writing $f(t)=(x(t),y(t))$ as a vector
$\vec{p}(t)=\langle x(t),y(t)\rangle$ and interpreting its magnitude.}
\end{skillbox}

\begin{skillbox}[Lesson]{sky}
\begin{multicols}{2}

\textbf{Step 1 --- Position as a Vector-Valued Function}

Any parametric function $f(t)=(x(t),y(t))$ can be rewritten as
\[
  \vec{p}(t) = \langle x(t),\; y(t)\rangle
             = x(t)\,\hat{\imath} + y(t)\,\hat{\jmath}.
\]
The \textbf{magnitude} equals the particle's distance from the origin:
\[
  \|\vec{p}(t)\| = \sqrt{[x(t)]^2 + [y(t)]^2}.
\]

\columnbreak

\textbf{Step 2 --- Velocity Vector and Speed}

At time $t$, the velocity vector uses the component rates of change:
\[
  \vec{v}(t) = \langle x'(t),\; y'(t)\rangle.
\]
\begin{itemize}
  \item $x'(t)>0$: moving right; $x'(t)<0$: left.
  \item $y'(t)>0$: moving up; $y'(t)<0$: down.
  \item $\|\vec{v}(t)\|=\sqrt{[x'(t)]^2+[y'(t)]^2}$\quad is the \textbf{speed}.
\end{itemize}

\end{multicols}

\medskip
\textbf{Step 3 --- Worked Example: Drone}\\[3pt]
Drone: $x(t)=3t$, $y(t)=t^{2}-4$.

\begin{enumerate}
  \item \textbf{Position vector:} $\vec{p}(t)=\langle 3t,\;t^{2}-4\rangle$.
  \item \textbf{Evaluate at $t=2$:} $\vec{p}(2)=\langle 6,\;0\rangle$.
        Distance from origin: $\|\vec{p}(2)\|=\sqrt{36+0}=6$ m.
  \item \textbf{Velocity at $t=2$:} With $x'(t)=3$ and $y'(t)=2t$,
        $\vec{v}(2)=\langle 3,\;4\rangle$.
        Speed: $\|\vec{v}(2)\|=\sqrt{9+16}=5$ m/s.
  \item \textbf{Interpret:} $x'(2)=3>0$ (moving right); $y'(2)=4>0$ (moving up).
\end{enumerate}
\end{skillbox}

\begin{skillbox}[Active Monitoring]{redbox}
\begin{itemize}
  \item Circulate while students convert a new parametric function to vector form;
        watch for confusion between $\vec{p}(t)$ (position) and $\vec{v}(t)$ (velocity).
  \item Check that students square \emph{both} components when computing magnitude
        (common error: computing $\sqrt{x(t)+y(t)}$ instead of
        $\sqrt{[x(t)]^2+[y(t)]^2}$).
  \item Cold-call prompt: ``What does the sign of $y'(t)$ tell us about the drone's
        altitude?  What does $\|\vec{v}(t)\|$ represent?''
\end{itemize}
\end{skillbox}

\begin{skillbox}[Group Work \& Differentiation]{redbox}
Students work in tiered groups using the tiered activity sheet.

\medskip
\begin{tcolorbox}[colback=white, colframe=black!40,
    title=\textbf{Tier R --- Remediate},
    fonttitle=\bfseries, arc=2mm, left=3mm, right=3mm, top=1mm, bottom=1mm]
Given $\vec{p}(t)=\langle 3t,\,2t\rangle$, evaluate at $t=0,1,2$ and find
$\|\vec{p}(2)\|$.  Focus: mechanical conversion and magnitude formula.
\end{tcolorbox}

\smallskip
\begin{tcolorbox}[colback=white, colframe=black!40,
    title=\textbf{Tier A --- Apply},
    fonttitle=\bfseries, arc=2mm, left=3mm, right=3mm, top=1mm, bottom=1mm]
Given $f(t)=(t^{2}-4,\,3t)$: write as a position vector, find when the
particle is on the $y$-axis, and find the speed at $t=1$.
\end{tcolorbox}

\smallskip
\begin{tcolorbox}[colback=white, colframe=black!40,
    title=\textbf{Tier E --- Extend},
    fonttitle=\bfseries, arc=2mm, left=3mm, right=3mm, top=1mm, bottom=1mm]
Show $\|\vec{p}(t)\|=2$ for all $t$ when $\vec{p}(t)=\langle 2\cos t,\,2\sin t\rangle$.
Interpret $\vec{v}(t)=\langle -2\sin t,\,2\cos t\rangle$ in terms of direction of
motion at $t=0$.
\end{tcolorbox}
\end{skillbox}

\begin{skillbox}[Individual Work \& Assessment]{redbox}
\textbf{Exit ticket} (individual, collected):

Given $\vec{p}(t)=\langle t^{2}-1,\,2t+3\rangle$:
\begin{enumerate}
  \item Find $\vec{p}(2)$ and explain what it represents in context.
  \item Compute $\|\vec{p}(2)\|$ to the nearest tenth.
  \item With $\vec{v}(t)=\langle 2t,\,2\rangle$, find the speed at $t=1$.
\end{enumerate}
\textit{Students who struggle with Question 2 likely need additional practice with
the magnitude formula before Lesson 4.10.}
\end{skillbox}

\begin{skillbox}[Reinforcement \& Extension]{goldbox}
\textbf{Homework:} Five problems — rewriting a parametric function as a
vector-valued function; finding position magnitude at a given $t$;
interpreting velocity vector signs; computing speed; AP-style MC on maximum
distance from the origin for $\vec{p}(t)=\langle 4\cos t,\,3\sin t\rangle$.

\medskip
\textbf{Preview Lesson 4.10 --- Matrices:} We used vectors to encode 2-D
position. Matrices generalize vectors and allow systematic transformation of
position and velocity data --- foundational for computer graphics, data
science, and physics.
\end{skillbox}

\end{document}
